% from ega4-errata-addenda-list3-fr.tex
\oldpage[IV]{358}
quotient de $X$ par la relation d'équivalence
$\overline{\{x\}}=\overline{\{x'\}}$. Il résulte aussitôt de cette
définition que tout ouvert de $X$ est saturé pour $R$, donc la topologie
de $X$ est l'image réciproque de celle de $Y$ par l'application
canonique $\pi:X\to X/R$. Un ouvert dense dans $X$ est de la forme
$\pi^{-1}(U)$, où $U$ est un ouvert dense dans $Y$, donc un fermé rare
$Z$ dans $X$ est de la forme $\pi^{-1}(S)$, où $S=\pi(Z)$ est fermé
rare dans $Y$, et il revient au même de dire que $X-Z$ est connexe ou
que $Y-S$ est connexe. De même $\pi(T_x)$ est l'ensemble des
générisations de $\pi(x)$, et il est équivalent de dire que
$T_x-\{x\}$ est connexe ou que $\pi(T_x)-\{\pi(x)\}$ l'est. Enfin, par
construction $Y$ est un espace de Kolmogoroff localement noethérien,
dont toute partie fermée admet donc un seul point générique. On voit
alors qu'il revient au même de prouver le lemme pour $X$ ou pour $Y$,
et on peut donc supposer que toute partie fermée irréductible de $X$ a
un point générique unique.
\end{quote}

\paragraph{(Err\textsubscript{IV}, 45).}
Dans (IV, 15.6.1), ligne 18 du bas de la p.~236, supprimer « localement
noethérien » (hypothèse non utilisée dans la démonstration).

\paragraph{(Err\textsubscript{IV}, 46).}
Dans (IV, 15.6.3), ligne 11 de la p.~237, au lieu de « supposons que
pour tout $y\in Y$ », lire : « supposons que $f$ soit localement de
présentation finie et que pour tout $y\in Y$ ». Remplacer la
démonstration (lignes 14 à 24 de la p.~237) par la suivante :
\begin{quote}
En vertu de {\hyperref[canon-index.ega4.c04.prop.1-10-3]{(1.10.3)}}, il suffit de
prouver que pour tout $x\in X_y^0$ et toute générisation $y'$ de $y$,
il existe une générisation $x'$ de $x$ telle que $f(x')=y'$. Or il
suffit de prendre pour $x'$ le point générique de $X_{y'}^0$, en vertu
de {\hyperref[canon-index.ega4.c04.cor.15-6-2]{(15.6.2)}}.
\end{quote}

\paragraph{(Err\textsubscript{IV}, 47).}
Dans (IV, 15.6.7), ligne 2 de la p.~241, après « noethérien », ajouter :
« ou si la condition $\beta$) de
({\hyperref[canon-index.ega4.c04.prop.15-6-6]{15.6.6}}, (iii)) est vérifiée ».
Après la ligne 9 de la p.~241, ajouter :
\begin{quote}
Lorsque la condition $\beta$) de
({\hyperref[canon-index.ega4.c04.prop.15-6-6]{15.6.6}}, (iii)) est satisfaite, le
fait que a) entraîne b') résulte de
({\hyperref[canon-index.ega4.c04.prop.15-6-6]{15.6.6}}, (ii)), et le fait que b')
entraîne a) résulte de
({\hyperref[canon-index.ega4.c04.prop.15-6-6]{15.6.6}}, (iii)).
\end{quote}

\paragraph{(Err\textsubscript{IV}, 48).}
Dans (IV, 20.1.5), ligne 13 de la p.~228, (IV, 20.1.6), ligne 17 de la
p.~228, (IV, 20.1.7), ligne 13 du bas de la p.~228, supprimer
« strictement ».

\paragraph{(Err\textsubscript{IV}, 49).}
Dans (IV, 20.2.13), ajouter après la ligne 6 de la p.~236 : En
particulier si $X$ est réduit et n'a qu'un nombre fini de composantes
irréductibles, la notion de $\sh{O}_X$-Module sans torsion définie dans
({\hyperlink{ega4.c04.statement.20-1-5}{IV, 20.1.5}}) coïncide avec
celle de (I, 7.4.1).

\paragraph{(Err\textsubscript{IV}, 50).}
Dans (IV, 20.2.14), lignes 18 et 25 de la p.~236, supprimer
« strictement ».

\paragraph{(Err\textsubscript{IV}, 51).}
Dans (IV, 20.6.2), ligne 13 du bas de la p.~252, supprimer
« strictement ».

\paragraph{(Err\textsubscript{IV}, 52).}
Dans (IV, 20.6.4), lignes 23 et 22 de la p.~253, supprimer
« strictement ».

\paragraph{(Err\textsubscript{IV}, 53).}
Dans (IV, 21.4), après la ligne 23 de la p.~267, ajouter :

La proposition suivante et ses corollaires nous ont été signalés par
M.~Raynaud.

\begin{proposition}[21.4.9]
\EGAReaderTarget{ega4.c04.prop.21-4-9}{IV.21.4.9}
Soient $X$, $Y$ deux préschémas localement noethériens,
$f:X\to Y$ un morphisme fidèlement plat; on suppose en outre, soit que
$f$ est quasi-compact, soit que $f$ est ouvert. Pour qu'un diviseur
$D'$ sur $X$ soit de la forme $f^*(D)$, où $D$ est un diviseur sur
$Y$, il faut et il suffit que pour tout $y\in Y$ tel que
$\operatorname{prof}(\sh{O}_{Y,y})\leq1$, en posant
\[
 Y_1=\operatorname{Spec}(\sh{O}_{Y,y}),\qquad
 X_1=X\times_YY_1,\qquad
 f_1=f_{(Y_1)}:X_1\to Y_1,
\]
l'image réciproque $D'_1$ de $D'$ sur $X_1$ soit de la forme
$f_1^*(D_1)$, où $D_1\in\operatorname{Div}(Y_1)$.
\end{proposition}

